\labsection{8}{Современная безградиентная оптимизация: CMA-ES, байесовская оптимизация и NSGA-II}

\labpart{Цель работы}
Изучить эволюционную стратегию с адаптацией ковариационной матрицы, байесовскую оптимизацию и
многокритериальный генетический алгоритм NSGA-II; сравнить их с генетическими алгоритмами ЛР~№5--6 и
применить к двум практическим задачам машинного обучения~--- подбору гиперпараметров и отбору признаков.

\labpart{Теоретические сведения}
Подраздел~\ref{sec:real} и раздел~\ref{sec:modern} теоретической части: CMA-ES
(\ref{eq:cma-sample})--(\ref{eq:cma-sigma}) и перезапуски IPOP, байесовская оптимизация и TPE, практика
подбора гиперпараметров, доминирование по Парето, NSGA-II, операторы SBX и полиномиальной мутации,
показатели HV и IGD.

\labpart{Постановка задачи}
\emph{Часть А.} Реализовать CMA-ES и IPOP-CMA-ES на NumPy и найти минимум функции варианта
(таблица~\ref{tab:l6-var}); сравнить со случайным поиском и вещественным ГА ЛР~№6 при одинаковом бюджете
вычислений функции; проверить реализацию по библиотеке \code{cma}.

\emph{Часть Б.} Подобрать два гиперпараметра модели варианта (таблица~\ref{tab:l8-var}) при бюджете 30
обучений: перебором по сетке, случайным поиском, своим вещественным ГА, Optuna TPE и Optuna CMA-ES;
целевая функция~--- точность на 5-блочной кросс-валидации обучающей части (80\,\%), итог проверяется на
отложенных 20\,\%.

\emph{Часть В.} Реализовать NSGA-II, проверить её на задаче ZDT1 (30 переменных) против
\code{pymoo} по показателям HV и IGD. Выполнить многокритериальный отбор признаков для датасета варианта:
минимизировать ошибку кросс-валидации логистической регрессии и число признаков (двоичная хромосома).

\begin{table}[!htb]
\caption{Варианты заданий к лабораторной работе №8 (функция~--- по таблице~\ref{tab:l6-var}, датасеты из
\code{sklearn.datasets}, модели~--- по таблице~\ref{tab:l8-models})}
\label{tab:l8-var}
\centering\small
\begin{tabular}{cccl@{\qquad}cccl}
\toprule
№ & Функция & Модель & Датасет & № & Функция & Модель & Датасет \\
\midrule
1 & 1 & SVC & \code{wine}           & 5 & 5 & HGB & \code{breast\_cancer} \\
2 & 2 & KNN & \code{digits}         & 6 & 6 & SVC & \code{digits} \\
3 & 3 & SVC & \code{breast\_cancer} & 7 & 7 & KNN & \code{breast\_cancer} \\
4 & 4 & RF  & \code{wine}           & 8 & 8 & HGB & \code{wine} \\
\bottomrule
\end{tabular}
\end{table}

\begin{table}[!htb]
\caption{Модели и гиперпараметры части Б: классы из \code{sklearn.svm}, \code{sklearn.neighbors} и
\code{sklearn.ensemble}; целочисленные параметры при поиске округляются}
\label{tab:l8-models}
\centering\small
\begin{tabularx}{\textwidth}{@{}ll>{\raggedright\arraybackslash}X@{}}
\toprule
Код & Класс & Гиперпараметры и диапазоны \\
\midrule
SVC & \code{SVC} (RBF-ядро) & $\lg C \in [-2; 4]$, $\lg\gamma \in [-6; 0]$ \\
KNN & \code{KNeighborsClassifier} & \code{n\_neighbors} $\in [1; 50]$, \code{p} $\in [1; 3]$ \\
RF  & \code{RandomForestClassifier} & \code{max\_depth} $\in [1; 20]$, \code{min\_samples\_leaf} $\in [1; 20]$ \\
HGB & \code{HistGradientBoostingClassifier} & $\lg$\,\code{learning\_rate} $\in [-3; 0]$, \code{max\_leaf\_nodes} $\in [2; 64]$ \\
\bottomrule
\end{tabularx}
\end{table}

\labpart{Требования к программе}
CMA-ES и NSGA-II реализуются самостоятельно; библиотеки \code{cma}, Optuna и \code{pymoo} используются
для сравнения. Отображаются: эллипсы распределения CMA-ES на линиях уровня функции; медианные кривые
<<лучшее значение от числа вычислений>> для всех методов; карта точности модели на плоскости двух
гиперпараметров с точками, выбранными методами поиска; найденный и истинный фронты Парето.

\labpart{Требования к интерфейсу и анимации}
\begin{itemize}
\item \emph{Графическое поле.} Часть~А~--- линии уровня функции, точки текущего поколения и эллипс распределения CMA-ES; часть~Б~--- карта точности на плоскости двух гиперпараметров с пробами метода; часть~В~--- плоскость критериев с популяцией и истинным фронтом. Рядом~--- график лучшего значения от числа вычислений.
\item \emph{Интерактивность.} Элементы управления: часть работы и метод поиска, размер популяции $\lambda$, начальный шаг $\sigma_0$, бюджет вычислений, seed.
\item \emph{Мультипликация.} Эллипс распределения смещается, поворачивается и сжимается от поколения к поколению; пробы гиперпараметров появляются на карте по одной; популяция NSGA-II приближается к истинному фронту.
\item \emph{Подсветка событий.} В части~А выделяются $\mu$ лучших точек поколения, по которым обновляется среднее, и отмечаются перезапуски IPOP; в части~Б~--- новая и лучшая пробы; в части~В недоминируемые решения отличаются цветом от остальных.
\end{itemize}

\labpart{Порядок выполнения работы}
\begin{steps}
\item Реализовать CMA-ES (приложение~\ref{app:listings}), проверить на функции сферы в 10-мерном
пространстве (ожидаемое $f < 10^{-20}$) и на функции Розенброка; построить эллипсы распределения.
\item Сравнить на функции варианта случайный поиск, вещественный ГА, CMA-ES, IPOP-CMA-ES и \code{cma}
(20 запусков, одинаковый бюджет), найти долю успешных запусков и число вычислений до успеха.
\item Для модели варианта построить карту точности на сетке $31 \times 31$ и выполнить поиск всеми
методами (10 запусков с разными seed); сравнить скорость роста лучшей точности и точность на тесте.
\item Реализовать NSGA-II и операторы SBX и полиномиальной мутации, сравнить с \code{pymoo} на ZDT1.
\item Выполнить отбор признаков с помощью NSGA-II (кэшируя оценки уже встречавшихся наборов) и выбрать
компромиссный набор признаков по фронту Парето с учётом точности на тесте.
\end{steps}

\labpart{Пример выполнения}
\emph{Часть А.} Вариант~3: функция Шаффера F7, бюджет 9060 вычислений (как у вещественного ГА ЛР~№6),
20 запусков. Своя CMA-ES совпадает по поведению с библиотечной (таблица~\ref{tab:l8-cma}), а на
функции Розенброка эллипс распределения вытягивается вдоль изогнутого оврага и сжимается у минимума
(рисунок~\ref{fig:l8-cmaes}).

\begin{table}[!htb]
\caption{Функция F7: лучшее $f$ за 9060 вычислений (20 запусков), успех~--- $f < 10^{-3}$}
\label{tab:l8-cma}
\centering\small
\begin{tabular}{lcccc}
\toprule
Метод & Медиана $f$ & Худшее $f$ & Успех, \% & Вычислений до успеха \\
\midrule
случайный поиск & 1,3 & 1,8 & 0 & --- \\
вещественный ГА (ЛР~№6) & $3,5 \cdot 10^{-9}$ & 0,015 & 90 & 3210 \\
CMA-ES (своя) & 0,010 & 10 & 15 & 714 \\
CMA-ES (\code{cma}) & 0,014 & 4,8 & 10 & --- \\
IPOP-CMA-ES (своя) & $5,6 \cdot 10^{-7}$ & 0,031 & 80 & 3456 \\
\bottomrule
\end{tabular}
\end{table}

\begin{figure}[!htb]
\centering
\includegraphics[width=\textwidth]{\figdir/lab8_cmaes.png}
\caption{Слева~--- эллипсы распределения CMA-ES на функции Розенброка (числа~--- номер поколения);
справа~--- медиана лучшего значения F7 по 20 запускам}
\label{fig:l8-cmaes}
\end{figure}

Одиночная CMA-ES в успешных запусках достигает $10^{-3}$ за 714 вычислений~--- в четыре раза быстрее
ГА, но в 17 из 20 запусков её распределение сжимается на одном из колец F7. Перезапуски с удвоением
популяции поднимают долю успехов до 80\,\%. На многоэкстремальной функции популяционный ГА остаётся
надёжнее, а CMA-ES выигрывает на гладких и овражных функциях.

\emph{Часть Б.} Вариант~3: \code{SVC} с RBF-ядром на \code{breast\_cancer}, $\lg C \in [-2; 4]$,
$\lg\gamma \in [-6; 0]$. Максимум точности на сетке $31 \times 31$~--- $0,987$. Результаты приведены в
таблице~\ref{tab:l8-hpo} и на рисунке~\ref{fig:l8-hpo}.

\begin{table}[!htb]
\caption{Подбор гиперпараметров \code{SVC}: 30 обучений, 10 запусков (сетка~--- один детерминированный); последний столбец~--- медиана числа обучений до CV $\ge 0,98$}
\label{tab:l8-hpo}
\centering\small
\begin{tabular}{lcccc}
\toprule
Метод & CV, медиана & CV, худший & Тест, медиана & До CV $\ge 0,98$ \\
\midrule
сетка $6 \times 5$ & 0,985 & 0,985 & 0,974 & 14 \\
случайный поиск & 0,981 & 0,976 & 0,974 & 14 (9 из 10) \\
ГА ($N = 6$, 5 поколений) & 0,982 & 0,980 & 0,974 & 10,5 \\
Optuna TPE & 0,984 & 0,980 & 0,974 & 7 \\
Optuna CMA-ES & 0,982 & 0,982 & 0,974 & 6 \\
\bottomrule
\end{tabular}
\end{table}

\begin{figure}[!htb]
\centering
\includegraphics[width=\textwidth]{\figdir/lab8_hpo.png}
\caption{Слева~--- точность кросс-валидации на плоскости гиперпараметров и 30 проб TPE и случайного
поиска в одном запуске; справа~--- медиана лучшей найденной точности}
\label{fig:l8-hpo}
\end{figure}

Методы с моделью поверхности (TPE, CMA-ES) выходят на точность $0,98$ за 6--7 обучений, случайный поиск и
сетка~--- за 14. TPE сгущает пробы в диагональной <<долине>> высокой точности. Итоговая точность на
тесте у всех методов одинакова ($0,974$): различия в CV на уровне $0,003$ меньше разброса оценки, и
выбор максимума CV немного переобучается на валидацию. Практический вывод: экономия проявляется в
числе дорогих обучений, а не в итоговом качестве простой модели.

\emph{Часть В.} На ZDT1 (100 особей, 200 поколений, 5 запусков) своя NSGA-II получила IGD $= 0,0044$ и
HV $= 0,8707$, \code{pymoo}~--- $0,0051$ и $0,8682$ (HV истинного фронта $0,8714$), за 0,28 и 0,36~с.
Первая версия реализации без обмена переменными между потомками SBX давала IGD $= 0,41$: без обмена
кроссинговер не комбинирует переменные разных родителей~--- типичная ошибка, которую ловит сравнение с
библиотекой. Отбор признаков на \code{breast\_cancer} (40 особей, 25 поколений, 813 различных наборов,
5,5~с) дал фронт, показанный на рисунке~\ref{fig:l8-pareto} и в таблице~\ref{tab:l8-fs}.

\begin{figure}[!htb]
\centering
\includegraphics[width=\textwidth]{\figdir/lab8_pareto.png}
\caption{Слева~--- фронты задачи ZDT1; справа~--- фронт Парето отбора признаков и точность на тесте}
\label{fig:l8-pareto}
\end{figure}

\begin{table}[!htb]
\caption{Фронт Парето отбора признаков: логистическая регрессия на \code{breast\_cancer}}
\label{tab:l8-fs}
\centering\small
\begin{tabular}{lcccccccc}
\toprule
Признаков & 1 & 2 & 3 & 4 & 5 & 6 & 9 & все 30 \\
\midrule
Точность CV   & 0,921 & 0,958 & 0,971 & 0,974 & 0,978 & 0,980 & 0,982 & 0,978 \\
Точность тест & 0,904 & 0,956 & 0,974 & 0,974 & 0,965 & 0,956 & 0,965 & 0,982 \\
\bottomrule
\end{tabular}
\end{table}

Трёх--четырёх признаков достаточно для точности 0,97, пяти~--- для точности кросс-валидации модели на
всех 30 признаках. Дальнейший рост CV не подтверждается на тесте: при отборе по той же валидации, по
которой строится фронт, признаки подгоняются под шум. Компромисс выбирают в <<колене>> фронта~--- здесь
3--4 признака.

\labpart{Контрольные вопросы}
\begin{questions}
\item В каких задачах нужны безградиентные методы оптимизации? Когда градиентные методы предпочтительнее?
\item Опишите правило одной пятой для (1+1)-ЭС.
\item Запишите шаг CMA-ES: порождение точек, обновление среднего, ковариационной матрицы и шага. Зачем
нужны пути эволюции?
\item Почему CMA-ES инвариантна к монотонным преобразованиям функции и к поворотам координат? Почему
одиночная CMA-ES застревает на многоэкстремальных функциях и как это исправляет IPOP?
\item Что такое суррогатная модель и функция приобретения? Запишите ожидаемое улучшение. Как работает TPE?
\item Почему случайный поиск гиперпараметров эффективнее сетки? Почему лучшее значение кросс-валидации
оптимистично смещено?
\item Дайте определения доминирования и фронта Парето. Почему взвешенная сумма критериев не находит
невыпуклые части фронта?
\item Опишите недоминируемую сортировку, расстояние скученности и отбор в NSGA-II.
\item Что измеряют гиперобъём и IGD?
\end{questions}
